\documentclass[exercice]{thedocument}%
\usepackage{enumitem}%
\startdatakeys
\source{}
\BO{2026}
\cycle{4}
\competencesDuSocle{CA,MO}
\objectifsApprentissage{C4A1ca,C4A1cb,C4C1ba,C4C1bb,C4C1bc}
\automatismes{C4A1cb,C4A1cc,C4C1ba,C4C1bb,C4C1bc}
\enddatakeys
\begin{document}%
Voici quatre affirmations. Pour chacune d'entre elles, dire si elle est vraie
ou fausse et {\bfseries justifier}.
\begin{enumerate}[label=\arabic*.]
\item {\bfseries Affirmation 1 : }$\dfrac35+\dfrac12=\dfrac{3+1}{5+2}$
\item On considère la fonction $f:x\mapsto5-3x$.\par
{\bfseries Affirmation 2 : }l'image de $-1$ par $f$ est -2.\vskip5pt
\item On considère deux expériences aléatoires :
\begin{itemize}[label=$-$]
    \item {\itshape expérience 1 : }choisir au hasard un nombre entier compris entre
    1 et 11 (1 et 11 inclus).
    \item {\itshape expérience 2 : }lancer un dé équilibré à six faces numérotées
    de 1 à 6 et annoncer le nombre qui apparaît sur la face du dessus.
\end{itemize}
{\bfseries Affirmation 3 : }il est plus probable de choisir un nombre premier dans
l'expérience 1 que d'obtenir un nombre pair dans l'expérience 2.
\item {\bfseries Affirmation 4 : }pour tout nombre $x$,$$(2x+1)^2=(2x+3)(2x-1)$$
\end{enumerate}
\end{document}
